\section{The exact words of twenty-one results}
\label{sec:coq-statements}
%===========================================================================

This book states its theorems in its own words, and Section~\ref{sec:methodology} says why that's a
risk: the checker confirms a proof matches its statement. Whether my sentence matches the
statement is a question the checker never asks. So here are the statements themselves for twenty-one of the results the argument
leans on hardest. Each is copied from its file with the comments taken out and a few long lines
broken to fit the page, which Coq doesn't notice, and nothing else changed, and each comes with its file, a plain reading, and a line on whether the book's
printed version says exactly that, more, or less. Each source block has a plain reading beside it. Compare the definitions, the premises,
the quantifiers and the conclusion with the printed theorem. Checking the translation
against the source also takes enough of its notation to follow those parts; matching
the list of premises alone is not the whole test. These are the statements this book was checked against; if a
file differs later, the book means the one printed here.

\subsection{The definitions they rest on}
\label{sec:coq-definitions}

A statement is only as plain as the definitions inside it, and the readings below describe several
of those definitions in words. Here are the ones they lean on hardest, copied the same way: the
small machine's step, the meaning of the universal reading checker, the
simulations that nesting is about, the meaning of the CHSH property, and the universal base with
the interface that carries it.

{\raggedright The step of the small machine, with the checks it runs and \code{untouched},
minimal/EarnedCore.v. \texttt{prop}, \texttt{eval}, \texttt{fact}, \texttt{fact\_eqb}, \texttt{cost}
and \texttt{run} are as in Definitions~\ref{def:u34-props}, \ref{def:u34-state} and~\ref{def:u34-step}.\par}

{\scriptsize
\begin{verbatim}
Inductive instr : Type :=
| INC (c : ctr)
| DEC (c : ctr) (j : nat)
| HALT
| CHECK (p : prop) (c : ctr)
| COMMIT (p : prop) (c : ctr)
| CERTIFY.

Record core : Type := mkcore {
  ca : nat; cb : nat;
  va : nat; vb : nat;
  pc : nat;
  facts : list fact;
  chan : option fact;
  err : bool
}.

Definition val (k : core) (c : ctr) : nat :=
  match c with CA => ca k | CB => cb k end.

Definition ver (k : core) (c : ctr) : nat :=
  match c with CA => va k | CB => vb k end.

Definition claim (k : core) (p : prop) (c : ctr) : fact := mkfact p c (ver k c).

Definition write (k : core) (c : ctr) (n j : nat) : core :=
  match c with
  | CA => mkcore n (cb k) (S (va k)) (vb k) j (facts k) (chan k) (err k)
  | CB => mkcore (ca k) n (va k) (S (vb k)) j (facts k) (chan k) (err k)
  end.

Definition goto (k : core) (j : nat) : core :=
  mkcore (ca k) (cb k) (va k) (vb k) j (facts k) (chan k) (err k).

Definition trap (k : core) : core :=
  mkcore (ca k) (cb k) (va k) (vb k) (pc k) (facts k) (chan k) true.

Definition record_fact (k : core) (f : fact) : core :=
  mkcore (ca k) (cb k) (va k) (vb k) (S (pc k)) (f :: facts k) (chan k) (err k).

Definition commit_to (k : core) (f : fact) : core :=
  mkcore (ca k) (cb k) (va k) (vb k) (S (pc k)) (facts k) (Some f) (err k).

Definition fact_cap : nat := 16.

Definition check_ok (k : core) (p : prop) (c : ctr) : bool :=
  negb (err k) && eval p (val k c) && Nat.ltb (length (facts k)) fact_cap.

Definition commit_ok (k : core) (p : prop) (c : ctr) : bool :=
  negb (err k) && existsb (fact_eqb (claim k p c)) (facts k).

Definition certify_ok (k : core) : bool :=
  negb (err k) && match chan k with Some _ => true | None => false end.

Definition cexec (k : core) (i : instr) : core :=
  if err k then k else
  match i with
  | INC c => write k c (S (val k c)) (S (pc k))
  | DEC c j =>
      match val k c with
      | 0 => goto k (S (pc k))
      | S n => write k c n j
      end
  | HALT => k
  | CHECK p c => if check_ok k p c then record_fact k (claim k p c) else trap k
  | COMMIT p c => if commit_ok k p c then commit_to k (claim k p c) else trap k
  | CERTIFY => if certify_ok k then goto k (S (pc k)) else trap k
  end.

Definition fires (k : core) (i : instr) : bool :=
  match i with CERTIFY => certify_ok k | _ => false end.

Record state : Type := mkst { core_of : core; mu : nat; cert : bool }.

Definition exec (s : state) (i : instr) : state :=
  mkst (cexec (core_of s) i) (mu s + cost i) (cert s || fires (core_of s) i).

Definition untouched (s : state) (mid : list instr) (c : ctr) : Prop :=
  forall t1 i t2, mid = t1 ++ i :: t2 ->
    let k := core_of (run t1 s) in
    ver (cexec k i) c = ver k c /\ val (cexec k i) c = val k c.
\end{verbatim}}

In plain words: a core holds the two counters, their versions, the line number, the fact table,
the channel and the trap bit. \code{claim k p c} is the fact made of \texttt{p}, \texttt{c} and
\texttt{c}'s version now. CHECK passes when the trap is down, the checker says yes of the counter's
value, and the table holds fewer than 16 facts, and it adds the claim to the table. COMMIT passes
when the trap is down and that exact claim, version and all, is in the table, and it puts the
claim in the channel in place of whatever was there. CERTIFY takes no operand: it passes when the
trap is down and the channel holds a fact, any fact. A failing one sets the trap and changes nothing
else, and a core with the trap up stays as it is. \code{exec} runs the core's step, adds the
instruction's cost to the ledger, and raises the flag when a passing CERTIFY runs.
\code{untouched s mid c} says each instruction of \texttt{mid}, at the point where it runs from
\texttt{s}, leaves \texttt{c}'s value and version as they were.
Definitions~\ref{def:u34-step} and~\ref{def:u34-untouched} say exactly this.

{\raggedright \code{cg\_uholds}, the meaning of the universal reading checker, with the claims it
reads and the routine's start, coq/kernel/foundation/CompilerChecker.v.\par}

{\scriptsize
\begin{verbatim}
Inductive cg_uprop : Type :=
| UBase (q : G.cprop)
| URun (r : nat).

Definition cg_env_chk (m v : nat) : env nat nat :=
  fun x => if Nat.ltb x m then cg_expo (qs x) v else 0.

Definition cg_uholds (p : cg_uprop) (v : nat) : Prop :=
  match p with
  | UBase q => G.cholds q v
  | URun r =>
      exists ig R xS xB xT m k st,
        cg_rdec r = Some (ig, R, xS, xB, xT, m) /\
        k <= cg_expo (qs xT) v /\
        sss_steps (mm_sss_env eq_nat_dec) (ig, R) k (ig, cg_env_chk m v) st /\
        out_code (fst st) (ig, R) /\
        get_env (snd st) xB = 1
  end.
\end{verbatim}}

In plain words: a claim is \code{UBase q}, a word of the counter language, or \code{URun r}, a
number naming a reading routine. A word means what it meant on the small machine. A routine claim
holds of a value \texttt{v} when \texttt{r} decodes to a routine (its first line, its code, and
three register numbers and a register count), and the routine, started at its first line with each
register below the count holding the exponent of that register's prime in \texttt{v}, reaches in
some number of steps no larger than the exponent of the fuel register's prime in \texttt{v} a state
outside its code with 1 in its answer register. \code{cg\_ueval\_iff}, printed below, says the
checker's yes is exactly this, and Section~\ref{sec:presented-universal} says the same.

{\raggedright \code{cmpz\_hom}, the simulations Theorem~\ref{thm:parts-nesting} is about, with the
objects they join, coq/kernel/foundation/AxCore.v, AxComplete.v and CzCat.v. \texttt{T.run} is
\texttt{run} of minimal/ThieleComplete.v, printed further down.\par}

{\scriptsize
\begin{verbatim}
Record BPre (A : Type) : Type := {
  bp_leb : A -> A -> bool;
  bp_refl : forall x, bp_leb x x = true;
  bp_trans : forall x y z, bp_leb x y = true -> bp_leb y z = true ->
    bp_leb x z = true
}.

Definition bp_le {A} (P : BPre A) (x y : A) : Prop := bp_leb A P x y = true.

Record amachine (A : Type) (P : BPre A) : Type := mk_am {
  am_state : Type;
  am_move : Type;
  am_step : am_state -> am_move -> am_state;
  am_cost : am_move -> nat;
  am_rec : am_state -> A
}.

Definition am_pt {A P} (AM : amachine A P) (r : am_state AM -> bool) : T.machine :=
  T.mk_machine (am_state AM) (am_move AM) (am_step AM) (am_cost AM) r.

Definition am_bare {A P} (AM : amachine A P) : T.machine :=
  am_pt AM (fun _ => false).

Definition am_run {A P} (AM : amachine A P) (tr : list (am_move AM))
    (s : am_state AM) : am_state AM :=
  T.run (am_bare AM) tr s.

Record cmpz_obj : Type := mk_obj {
  ob_A : Type;
  ob_P : BPre ob_A;
  ob_M : amachine ob_A ob_P
}.

Record cmpz_hom (X Y : cmpz_obj) : Type := mk_hom {
  h_state : am_state (ob_M X) -> am_state (ob_M Y);
  h_rec : ob_A X -> ob_A Y;
  h_mono : forall a b, bp_le (ob_P X) a b -> bp_le (ob_P Y) (h_rec a) (h_rec b);
  h_rec_nat : forall s, am_rec (ob_M Y) (h_state s) = h_rec (am_rec (ob_M X) s);
  h_sim : forall s m, exists tr : list (am_move (ob_M Y)),
            h_state (am_step (ob_M X) s m) = am_run (ob_M Y) tr (h_state s)
}.
\end{verbatim}}

In plain words: an object is a type of records with an order given by a yes-or-no test, and a
machine whose record lives in that type. A morphism from $X$ to $Y$ is a map on states and a map on
records such that the record map keeps the order, reading the record after mapping the state gives
the record map of the old record, and each move of $X$ is matched by some list of moves of $Y$ that
lands on the image of the new state. That last field asks for some list and names none, which is why
Theorem~\ref{thm:parts-nesting}(1) and the statement printed for it below speak of some host
trace.

{\raggedright \code{small\_chsh\_meaning}, the meaning of the property PCHSH, with the correlators
and the score it's read with, coq/kernel/quantum/SmallChshCheck.v, and the three definitions it
uses from TsirelsonFromAlgebra.v and NPAMomentMatrix.v. \code{npa\_to\_matrix} lays the ten entries
out as the five-by-five matrix $M(E)$ of Section~\ref{sec:chsh}.\par}

{\scriptsize
\begin{verbatim}
Definition CHSH_value (e00 e01 e10 e11 : R) : R :=
  e00 + e01 + e10 - e11.

Definition npa_psd (npa : NPAMomentMatrix) : Prop :=
  let M := nat_matrix_to_fin5 (npa_to_matrix npa) in
  symmetric5 M /\ PSD5 M.

Definition zero_marginal_npa (E00 E01 E10 E11 : R) : NPAMomentMatrix := {|
  npa_EA0 := 0;
  npa_EA1 := 0;
  npa_EB0 := 0;
  npa_EB1 := 0;
  npa_E00 := E00;
  npa_E01 := E01;
  npa_E10 := E10;
  npa_E11 := E11;
  npa_rho_AA := 0;
  npa_rho_BB := 0;
|}.

Definition small_chsh_corr (same diff : nat) : R :=
  if Nat.eqb (same + diff) 0 then 0
  else (INR same - INR diff) / INR (same + diff).

Definition small_chsh_e00 (t : small_chsh_tally) : R :=
  small_chsh_corr t.(small_chsh_same00) t.(small_chsh_diff00).

Definition small_chsh_e01 (t : small_chsh_tally) : R :=
  small_chsh_corr t.(small_chsh_same01) t.(small_chsh_diff01).

Definition small_chsh_e10 (t : small_chsh_tally) : R :=
  small_chsh_corr t.(small_chsh_same10) t.(small_chsh_diff10).

Definition small_chsh_e11 (t : small_chsh_tally) : R :=
  small_chsh_corr t.(small_chsh_same11) t.(small_chsh_diff11).

Definition small_chsh_score (t : small_chsh_tally) : R :=
  CHSH_value (small_chsh_e00 t) (small_chsh_e01 t)
             (small_chsh_e10 t) (small_chsh_e11 t).

Definition small_chsh_sampled (t : small_chsh_tally) : Prop :=
  (0 < t.(small_chsh_same00) + t.(small_chsh_diff00))%nat /\
  (0 < t.(small_chsh_same01) + t.(small_chsh_diff01))%nat /\
  (0 < t.(small_chsh_same10) + t.(small_chsh_diff10))%nat /\
  (0 < t.(small_chsh_same11) + t.(small_chsh_diff11))%nat.

Definition small_chsh_psd (t : small_chsh_tally) : Prop :=
  npa_psd (zero_marginal_npa (small_chsh_e00 t) (small_chsh_e01 t)
                             (small_chsh_e10 t) (small_chsh_e11 t)).

Definition small_chsh_meaning (t : small_chsh_tally) : Prop :=
  small_chsh_sampled t /\ small_chsh_psd t.
\end{verbatim}}

In plain words: the correlator of a question pair is the number of rounds that agreed minus the
number that disagreed, over the number of rounds, read as 0 when the pair had none. The score is the
CHSH combination of the four correlators. The meaning is that every question pair has at least one
round and the pinned matrix of the four correlators, with every one-sided average and both
same-side entries set to 0, is symmetric and positive semidefinite.
Definition~\ref{def:chsh-packed} says exactly this.

{\raggedright The universal base and the interface that carries it, minimal/ThieleComplete.v.
\texttt{machine} has the fields \texttt{m\_state}, \texttt{m\_move}, \texttt{m\_step},
\texttt{m\_cost} and \texttt{m\_record}.\par}

{\scriptsize
\begin{verbatim}
Inductive reg : Type := RA | RB.

Inductive cm_instr : Type := CINC (r : reg) | CDEC (r : reg) (j : nat).

Definition cm_conf : Type := (nat * (nat * nat))%type.

Definition cm_get (x : cm_conf) (r : reg) : nat :=
  match r with RA => fst (snd x) | RB => snd (snd x) end.

Definition cm_set (x : cm_conf) (r : reg) (n j : nat) : cm_conf :=
  match r with RA => (j, (n, snd (snd x))) | RB => (j, (fst (snd x), n)) end.

Definition cm_exec (i : cm_instr) (x : cm_conf) : cm_conf :=
  match i with
  | CINC r => cm_set x r (S (cm_get x r)) (S (fst x))
  | CDEC r j =>
      match cm_get x r with
      | 0 => (S (fst x), snd x)
      | S n => cm_set x r n j
      end
  end.

Record universal_base (M : machine) : Type := mk_ub {
  ub_window : m_state M -> cm_conf;
  ub_live : m_state M -> Prop;
  ub_compile : cm_instr -> m_move M;
  ub_load : nat -> nat -> m_state M;
  ub_load_window : forall a b, ub_window (ub_load a b) = (1, (a, b));
  ub_load_live : forall a b, ub_live (ub_load a b);
  ub_sim : forall s i, ub_live s ->
    ub_window (m_step M s (ub_compile i)) = cm_exec i (ub_window s) /\
    ub_live (m_step M s (ub_compile i))
}.

Inductive kind (C : Type) : Type :=
| KBase
| KCheck (c : C)
| KCommit (c : C)
| KCertify.

Record thiele_interface (M : machine) : Type := mk_ti {
  ti_base : universal_base M;
  ti_claim : Type;
  ti_kind : m_move M -> kind ti_claim;
  ti_meaning : ti_claim -> m_state M -> Prop;
  ti_check : m_state M -> ti_claim -> bool;
  ti_same : ti_claim -> m_state M -> m_state M -> Prop;
  ti_clean : m_state M -> Prop;
  ti_ledger : m_state M -> nat
}.
\end{verbatim}}

In plain words: a two-counter configuration is a line number and two registers, and
\code{cm\_exec} is what one instruction does to it, the same INC and DEC-or-fall-through as the
small machine's base. A universal base for a machine reads a configuration off each state (its
window), picks out the live states, gives a move for each instruction and a loaded state for each
start; each loaded state shows $(1, a, b)$ and is live, and from every live state the compiled move
changes the window exactly as \code{cm\_exec} says and leaves the state live. The interface lists
its parts and asks nothing of them: the base, the claims, a kind for each move, each claim's
meaning, the checker, ``unchanged,'' the clean states and the ledger. The demands are the four
clauses printed below under certified Thiele-complete. This is part
(a) of Definition~\ref{def:u34-thiele-complete} and the interface it describes.

\subsection{The universal floor}

{\raggedright \code{universal\_nfi\_any\_substrate}, coq/kernel/nfi/UniversalCertificationCost.v, with the record it quantifies over.\par}

{\footnotesize
\begin{verbatim}
Record CertificationSystem := mk_cert_system {
  cs_state : Type;
  cs_instr : Type;
  cs_step  : cs_state -> cs_instr -> cs_state;
  cs_cost  : cs_instr -> nat;
  cs_cert  : cs_state -> bool;
  cs_cert_costs :
    forall (s : cs_state) (i : cs_instr),
      cs_cert s = false ->
      cs_cert (cs_step s i) = true ->
      cs_cost i >= 1;
}.

Theorem universal_nfi_any_substrate :
  forall (CS : CertificationSystem)
         (trace : list (cs_instr CS))
         (s0 : cs_state CS),
    cs_cert CS s0 = false ->
    cs_cert CS (cs_run CS trace s0) = true ->
    cs_total_cost CS trace >= 1.
\end{verbatim}}

In plain words: a Thiele machine is a record of a state type, a move type, a step, a
cost and a yes-or-no reading, plus one more field, \code{cs\_cert\_costs}, which is the toll: a
move that takes the reading from false to true costs at least one. Since the toll is a field,
every system the theorem talks about has it, which is why the same file also states it as
\code{every\_toll\_system\_pays\_the\_floor}. The theorem: start at a state reading false, run
any list of moves (\code{cs\_run} applies them in order, \code{cs\_total\_cost} adds their
costs), end at a state reading true, and the costs add up to at least one.
Observation~\ref{thm:universal-floor} says exactly this; the rule that Definition~\ref{def:certsys}
calls the toll is the record's last field here.

\subsection{A permanent flip merges}

{\raggedright \code{permanent\_flip\_is\_not\_injective}, coq/kernel/nfi/PermanentCertification.v, inside a section that fixes the states, the moves, the step and the reading.\par}

{\footnotesize
\begin{verbatim}
Variables (S I : Type).
Variable step : S -> I -> S.
Variable cert : S -> bool.

Definition finite_states (all : list S) : Prop :=
  NoDup all /\ forall s, In s all.

Definition permanent : Prop :=
  forall s i, cert s = true -> cert (step s i) = true.

Definition step_injective (i : I) : Prop :=
  forall a b, step a i = step b i -> a = b.

Theorem permanent_flip_is_not_injective :
  forall all s i,
    finite_states all ->
    permanent ->
    cert s = false ->
    cert (step s i) = true ->
    ~ step_injective i.
\end{verbatim}}

In plain words: \texttt{finite\_states} $\mathit{all}$ says the list $\mathit{all}$ has no repeats and every
state is on it, so there are finitely many. \code{permanent} says once the reading is true,
every move keeps it true. \texttt{step\_injective i} says move $i$ never sends two different
states to the same one. The theorem: on a finite machine with a permanent reading, a move that
takes some state from false to true isn't injective, so it merges two states.
Theorem~\ref{thm:perm-flip-merges} says exactly this, with ``merges'' meaning ``isn't
injective''.

\subsection{The toll from finiteness, permanence, and merge pricing}

{\raggedright \code{a2\_from\_merging\_price\_and\_permanence}, the same file and section, with a cost on moves added.\par}

{\footnotesize
\begin{verbatim}
Variable cost : I -> nat.

Definition merging_steps_priced : Prop :=
  forall i, ~ step_injective i -> cost i >= 1.

Definition a2_holds : Prop :=
  forall s i, cert s = false -> cert (step s i) = true -> cost i >= 1.

Theorem a2_from_merging_price_and_permanence :
  forall all,
    finite_states all ->
    permanent ->
    merging_steps_priced ->
    a2_holds.
\end{verbatim}}

In plain words: \code{merging\_steps\_priced} says every move that isn't injective costs at
least one, and \code{a2\_holds} is the toll. The theorem: finite, permanent and merges priced
give the toll. Observation~\ref{thm:toll-derived} says this and one thing more, that the machine
with that cost is then a Thiele machine the universal floor applies to. That second half
isn't in this statement. It's the definition right after it in the file
(\code{certification\_system\_from\_merging\_price}), which fills the record's toll field with
this theorem, so the printed form is true, and it leans on two Coq items: this theorem and that
definition.

\subsection{No exact window price across a collision}

{\raggedright \code{shadow\_cannot\_price\_exactly}, coq/kernel/nfi/ShadowPricing.v, inside a section that fixes the states, moves, step, reading and a window \code{obs}.\par}

{\footnotesize
\begin{verbatim}
Variables (S I O : Type).
Variable step : S -> I -> S.
Variable cert : S -> bool.
Variable obs : S -> O.

Definition flips (s : S) (i : I) : bool :=
  negb (cert s) && cert (step s i).

Definition shadow_cost (price : O -> O -> nat) (s : S) (i : I) : nat :=
  price (obs s) (obs (step s i)).

Definition meets_floor (cost : S -> I -> nat) : Prop :=
  forall s i, flips s i = true -> cost s i >= 1.

Definition never_overcharges (cost : S -> I -> nat) : Prop :=
  forall s i, flips s i = false -> cost s i = 0.

Definition shadow_collision : Prop :=
  exists s1 i1 s2 i2,
    obs s1 = obs s2 /\
    obs (step s1 i1) = obs (step s2 i2) /\
    flips s1 i1 = true /\
    flips s2 i2 = false.

Theorem shadow_cannot_price_exactly :
  shadow_collision ->
  forall price,
    ~ (meets_floor (shadow_cost price) /\
       never_overcharges (shadow_cost price)).
\end{verbatim}}

In plain words: \texttt{flips s i} says move $i$ takes $s$ from false to true. A shadow price
is any function of what the window shows before and after a step. \code{meets\_floor} says
every flip costs at least one, \code{never\_overcharges} says every step that isn't a flip costs
nothing, and a \code{shadow\_collision} is two steps that look the same through the window,
before and after, one a flip and one not. The theorem: given a collision, no shadow price does
both. Observation~\ref{thm:shadow-price} bundles this with three neighbours from the same file:
that a price meeting the floor charges some non-flip at least one (\code{shadow\_floor\_overcharges}),
that a window showing the reading prices exactly
(\code{window\_showing\_reading\_prices\_exactly}), and that such a window has no collision
(\code{window\_showing\_reading\_has\_no\_collision}). This statement is the second of the four
clauses, and the book says exactly what it says.

\subsection{The record is a latch}

{\raggedright \code{record\_axis\_is\_latch\_holds}, coq/kernel/foundation/StructuralRecordAxis.v; the definitions come from StructuralCoreAnyBase.v and StructuralCoreCover.v.\par}

{\scriptsize
\begin{verbatim}
Definition record_permanent (M : RCM) : Prop :=
  forall s, rc_cert M s = true -> rc_cert M (rc_next M s) = true.

Definition computation_driven (M : RCM) (B : BaseMachine) (C : BaseCover M B)
  : Prop :=
  exists f : b_state B -> bool -> bool,
    forall m, rc_cert M (rc_next M m) = f (base_state M B C m) (rc_cert M m).

Definition HonestBaseExtension (M : RCM) (B : BaseMachine) (C : BaseCover M B)
  : Prop :=
  computation_driven M B C /\ ledger_carried M /\ rc_a2 M /\
  record_permanent M /\ reachable_record_write M.

Definition latch_next (B : BaseMachine) (h : b_state B -> bool)
    (x : b_state B * bool) : b_state B * bool :=
  (b_next B (fst x), orb (snd x) (h (fst x))).

Definition latch_factorization (M : RCM) (B : BaseMachine) (C : BaseCover M B)
    (h : b_state B -> bool) : Prop :=
  forall m,
    (base_state M B C (rc_next M m), rc_cert M (rc_next M m)) =
    latch_next B h (base_state M B C m, rc_cert M m).

Definition record_axis_is_latch : Prop :=
  forall M B C, HonestBaseExtension M B C ->
    exists h, latch_factorization M B C h.

Theorem record_axis_is_latch_holds : record_axis_is_latch.
\end{verbatim}}

In plain words: here a machine $M$ steps by itself (\code{rc\_next} gives the one next
state), carries a reading \code{rc\_cert} and a ledger, and sits over a base $B$ through a
cover \code{C}, a map down to the base's states that commutes with the step.
\code{computation\_driven} says the reading's next value is some function \code{f} of the base
state and the reading's current value. \code{HonestBaseExtension}, the book's base-driven record, means driven, a ledger that never falls, the
toll, a permanent reading, and some reachable step that writes it. \code{latch\_next} is the
latch of an event \code{h}: the base steps, and the record becomes ``already on, or \code{h}
holds now.'' The theorem: every base-driven record evolves exactly as the latch of some event.
Observation~\ref{thm:u11-latch} says a little more than this statement: it names the event,
$h(b) = f(b, \mathrm{off})$, where the statement only says one exists. The proof builds exactly
that event (its first move is \texttt{exists (fun b => f b false)}), so it shows the theorem's
stronger form too. And the proof uses two of the five
conditions, driven and permanent; the other three are carried along and never used. The form
with only those two as hypotheses is \code{record\_axis\_is\_latch\_r1\_r4}
(StructuralClockSync.v), and that is the form Observation~\ref{thm:u11-latch} states.

\subsection{A raised flag was checked, then committed, then certified}

{\raggedright \code{earned\_certification\_same\_claim}, minimal/EarnedCore.v.\par}

{\tiny
\begin{verbatim}
Theorem earned_certification_same_claim : forall s0 tr,
  clean_start s0 -> cert (run tr s0) = true ->
  exists pre1 p c mid1 mid2 post,
    tr = pre1 ++ CHECK p c :: mid1 ++ COMMIT p c :: mid2 ++ CERTIFY :: post /\
    check_ok (core_of (run pre1 s0)) p c = true /\
    commit_ok (core_of (run (pre1 ++ CHECK p c :: mid1) s0)) p c = true /\
    certify_ok (core_of (run (pre1 ++ CHECK p c :: mid1 ++ COMMIT p c :: mid2) s0))
      = true /\
    chan (core_of (run (pre1 ++ CHECK p c :: mid1 ++ COMMIT p c :: mid2) s0))
      = Some (claim (core_of (run (pre1 ++ CHECK p c :: mid1) s0)) p c) /\
    claim (core_of (run (pre1 ++ CHECK p c :: mid1) s0)) p c
      = mkfact p c (ver (core_of (run pre1 s0)) c) /\
    ver (core_of (run pre1 s0)) c = ver (core_of (run (pre1 ++ CHECK p c :: mid1) s0)) c /\
    untouched (run (pre1 ++ [CHECK p c]) s0) mid1 c.
\end{verbatim}}

In plain words: the clean-start premise on \texttt{s0} says the fact table and the channel are
empty and the flag is down, and \texttt{run tr s0} runs a list of instructions. The theorem: from a clean start,
if the flag is up at the end, the list splits as some start, CHECK \texttt{p c}, a middle, COMMIT
of the same \texttt{p c}, another middle, CERTIFY, and the rest; the CHECK, the COMMIT and the
CERTIFY each pass where they run; when the CERTIFY runs, the channel holds the claim the COMMIT
made, which is the fact (\texttt{p}, \texttt{c}, the version \texttt{c} had at the CHECK); the
counter's version at the CHECK equals its version at the COMMIT; and nothing in the first middle
touched that counter. CERTIFY takes no operand, and the channel clause is what ties it to the
claim. Theorem~\ref{thm:u34-cert-earned} says exactly this. The same file keeps the statement
without the channel clause as \code{earned\_certification\_provenance}.

\subsection{Checker soundness}

{\raggedright \code{checker\_soundness}, minimal/EarnedCore.v.\par}

{\scriptsize
\begin{verbatim}
Theorem checker_soundness : forall s0 tr f,
  clean_start s0 ->
  let k := core_of (run tr s0) in
  In f (facts k) -> f_ver f = ver k (f_ctr f) -> holds (f_prop f) (val k (f_ctr f)).
\end{verbatim}}

In plain words: from a clean start, after any list of instructions, every fact in the table
that's about its counter's current version is true of that counter's current value, where
\code{holds} is the property's meaning written without the checker.
Theorem~\ref{thm:u34-soundness} says exactly this in its first sentence; its ``in particular''
is the next corollary in the file, \code{committed\_claim\_holds}.

\subsection{No forging}

{\raggedright \code{no\_forging}, minimal/EarnedCore.v, with the two definitions it uses.\par}

{\scriptsize
\begin{verbatim}
Definition earned (s0 : state) (tr : list instr) (f : fact) : Prop :=
  exists pre mid,
    tr = pre ++ CHECK (f_prop f) (f_ctr f) :: mid /\
    check_ok (core_of (run pre s0)) (f_prop f) (f_ctr f) = true /\
    f = claim (core_of (run pre s0)) (f_prop f) (f_ctr f) /\
    (f_ver f = ver (core_of (run tr s0)) (f_ctr f) ->
     untouched (run (pre ++ [CHECK (f_prop f) (f_ctr f)]) s0) mid (f_ctr f)).

Definition no_forgery (s0 : state) (tr : list instr) : Prop :=
  forall f, In f (facts (core_of (run tr s0))) -> earned s0 tr f.

Theorem no_forging : forall s0 tr, clean_start s0 -> no_forgery s0 tr.
\end{verbatim}}

In plain words: a fact is \code{earned} when a passing CHECK of exactly that property and
counter wrote it somewhere in the list, it is exactly the claim that CHECK made, and, if it's
still about the counter's current version, nothing touched the counter since. \code{no\_forgery}
says every fact in the table at the end is earned. The theorem: from every clean start, every
list of instructions leaves no forgery. That's the last clause of
Theorem~\ref{thm:u34-no-forging}, with Definition~\ref{def:u34-earned} saying the same as
\code{earned}; the two one-step clauses before it are \code{facts\_step} (one instruction at a core) and \code{no\_forging\_step} (a trace with no forgery stays that way when one instruction is added).

\subsection{The minimum price of a certificate}

{\raggedright \code{certified\_run\_min\_cost}, minimal/EarnedCore.v.\par}

{\footnotesize
\begin{verbatim}
Theorem certified_run_min_cost : forall s0 tr,
  clean_start s0 -> cert (run tr s0) = true ->
  total_cost tr >= 3 /\ mu (run tr s0) >= mu s0 + 3.
\end{verbatim}}

In plain words: from a clean start, a list of instructions that ends with the flag up cost at
least three in total, and the ledger rose by at least three. Observation~\ref{thm:u34-min-price}
says exactly this in its first sentence; its form for stored programs is
\code{program\_certified\_min\_cost}.

\subsection{Three, reached by a check that could have failed}

{\raggedright \code{earned\_run\_check\_can\_fail} and \code{earned\_run\_refused\_forever}, minimal/EarnedCore.v, with the program they run.\par}

{\scriptsize
\begin{verbatim}
Definition earned_run : list instr := [CHECK PZero CA; COMMIT PZero CA; CERTIFY].

Theorem earned_run_check_can_fail : forall a b,
  (a = 0 ->
     cert (run_prog 4 earned_run (start a b)) = true /\
     mu (run_prog 4 earned_run (start a b)) = 3) /\
  (a <> 0 ->
     cert (run_prog 4 earned_run (start a b)) = false /\
     err (core_of (run_prog 4 earned_run (start a b))) = true).

Theorem earned_run_refused_forever : forall n a b,
  a <> 0 -> cert (run_prog n earned_run (start a b)) = false.
\end{verbatim}}

In plain words: \code{earned\_run} is CHECK ``$A$ is 0'', COMMIT the same, CERTIFY, and
\texttt{start a b} puts $a$ in $A$ and $b$ in $B$ with an empty table. The first theorem: run it
for four program steps; if $a$ is 0 the flag is up and the ledger is exactly 3, and if $a$
isn't 0 the flag is down and the trap is up. The second: if $a$ isn't 0, the flag is down after
any number of steps. So three is reached by a check that could have failed, and when it fails
nothing ever certifies. Theorem~\ref{thm:u34-tight} says this. The older witness in the same
file, \code{min\_cost\_tight}, reaches three with ``$A$ is at least 0'', a check that can't
fail; the step doesn't police that, and Section~\ref{sec:smallest} says so.

\subsection{Certified Thiele-complete, the four clauses}

{\raggedright \code{thiele\_complete}, minimal/ThieleComplete.v, with the clauses it is made of, and its form over a claim language fixed first, \code{thiele\_complete\_over}. \texttt{load a b} is the interface's loaded state for the start $(1, a, b)$, written inside a section that fixes \texttt{I}; \texttt{run M} runs a list of moves from a state.\par}

{\scriptsize
\begin{verbatim}
Definition universal_base_clause : Prop :=
  (forall i, ti_kind I (ub_compile (ti_base I) i) = KBase) /\
  (forall a b, ti_clean I (load a b)) /\
  (forall s m, ti_kind I m = KBase -> m_record M (m_step M s m) = m_record M s) /\
  (forall s m, m_record M s = true -> m_record M (m_step M s m) = true).

Definition earned_chain (s0 : m_state M) (tr : list (m_move M)) : Prop :=
  exists pre c chk mid1 cmt mid2 crt post,
    tr = pre ++ chk :: mid1 ++ cmt :: mid2 ++ crt :: post /\
    ti_kind I chk = KCheck c /\ ti_kind I cmt = KCommit c /\
      ti_kind I crt = KCertify /\
    ti_check I (run M pre s0) c = true /\
    (forall t1 t2, mid1 = t1 ++ t2 ->
       ti_same I c (run M pre s0) (run M (pre ++ chk :: t1) s0)) /\
    m_record M (run M (pre ++ chk :: mid1 ++ cmt :: mid2) s0) = false /\
    m_record M (run M (pre ++ chk :: mid1 ++ cmt :: mid2 ++ [crt]) s0) = true.

Definition earned_record_clause : Prop :=
  (forall s, ti_clean I s -> m_record M s = false) /\
  (forall s0 tr, ti_clean I s0 -> m_record M (run M tr s0) = true ->
    earned_chain s0 tr) /\
  (forall s c, ti_check I s c = true -> ti_meaning I c s) /\
  (forall c s s', ti_same I c s s' -> ti_meaning I c s -> ti_meaning I c s').

Definition record_move (m : m_move M) : nat :=
  match ti_kind I m with KBase => 0 | _ => 1 end.

Definition exact_toll_clause : Prop :=
  (forall m, m_cost M m = record_move m) /\
  (forall s m, ti_ledger I (m_step M s m) = ti_ledger I s + m_cost M m).

Definition non_vacuity_clause : Prop :=
  exists c chk cmt crt,
    ti_kind I chk = KCheck c /\ ti_kind I cmt = KCommit c /\
      ti_kind I crt = KCertify /\
    (forall a b, m_record M (run M [chk; cmt; crt] (load a b)) = true <->
                 ti_meaning I c (load a b)) /\
    (exists a b, ti_meaning I c (load a b)) /\
    (exists a b, ~ ti_meaning I c (load a b)).

Definition thiele_complete_with : Prop :=
  universal_base_clause /\ earned_record_clause /\ exact_toll_clause /\
    non_vacuity_clause.

Definition thiele_complete (M : machine) : Prop :=
  exists I : thiele_interface M, thiele_complete_with I.

Record claim_language (M : machine) : Type := mk_cl {
  cl_claim : Type;
  cl_eqb : cl_claim -> cl_claim -> bool;
  cl_eqb_spec : forall c d, cl_eqb c d = true <-> c = d;
  cl_meaning : cl_claim -> m_state M -> Prop;
  cl_same : cl_claim -> m_state M -> m_state M -> Prop;
  cl_same_keeps : forall c s s', cl_same c s s' ->
    cl_meaning c s -> cl_meaning c s'
}.

Record lang_interface (M : machine) (L : claim_language M) : Type := mk_lang {
  lang_base : universal_base M;
  lang_kind : m_move M -> kind (cl_claim L);
  lang_check : m_state M -> cl_claim L -> bool;
  lang_sound : forall s c, lang_check s c = true -> cl_meaning L c s;
  lang_clean : m_state M -> Prop;
  lang_ledger : m_state M -> nat
}.

Definition lang_ti {M : machine} {L : claim_language M}
  (J : lang_interface M L) : thiele_interface M :=
  mk_ti M (lang_base J) (cl_claim L) (lang_kind J) (cl_meaning L)
    (lang_check J) (cl_same L) (lang_clean J) (lang_ledger J).

Definition thiele_complete_over (M : machine) (L : claim_language M) : Prop :=
  exists J : lang_interface M L, thiele_complete_with (lang_ti J).

Theorem clock_not_thiele_complete : forall rd, ~ thiele_complete (clock rd).
\end{verbatim}}

In plain words: \texttt{M} is a machine (states, moves, step, cost, record) and \texttt{I} an
interface for it (a universal base, claims with meanings, a checker, ``unchanged,'' clean states,
a ledger, and a kind for each move). The first clause says the compiled two-counter moves are
base moves, loaded states are clean, base moves leave the record alone, and the record never comes
down; the universal base's own fields, printed in Section~\ref{sec:coq-definitions}, say the compiled moves act on the window
exactly as the two-counter instructions act. \code{earned\_chain} is a passing check of a claim,
then a commit of the same claim with the thing it's about unchanged at every point between, then
the certify that raises the record. The earned-record clause says clean states have the record
down, every run from a clean state that raises it contains that chain, a passing check means the
claim holds, and a claim that holds keeps holding while its thing is unchanged. The exact toll:
base moves cost 0, record moves 1, and the ledger adds each move's cost. Non-vacuity: some claim's
bare chain raises the record from a loaded state exactly when the claim holds there, and the claim
holds at one loaded state and fails at another. \code{thiele\_complete} asks for one interface
meeting all four. A \code{claim\_language} fixes the claims, an exact equality test on them, their
meanings, ``unchanged,'' and the proof that a meaning keeps holding while its thing is unchanged;
a \code{lang\_interface} over it supplies only the base, the kinds, the checker with its proof of
soundness against the language's meanings, the clean states and the ledger; \code{lang\_ti} makes
it a plain interface whose claims, meanings and ``unchanged'' are the language's; and
\code{thiele\_complete\_over} asks that interface to meet all four.
\texttt{clock} is the two-counter machine with every move costing 1 and any
reading of its configuration as the record. Definition~\ref{def:u34-thiele-complete} says
this, the language-first form included. The file also has a Coq predicate,
\code{weakly\_thiele\_complete} (not printed here), that asks only for the toll and a universal
base. Observation~\ref{thm:u34-clock} says the last line and more: the clock meets
\code{weakly\_thiele\_complete} (\code{clock\_weakly\_thiele\_complete}), a flag that switches on at the first yes
behaves the same way (\code{latch\_clock\_weakly\_thiele\_complete} and
\code{latch\_clock\_not\_thiele\_complete}), and no machine whose every move costs at least 1 is
certified Thiele-complete, which comes from \code{complete\_has\_free\_move}. Clause (c) fixes the unit at
one. Theorem~\ref{thm:u34-tc-unit} states the definition with every record move at one fixed price
$c \ge 1$ (\code{thiele\_complete\_at}, ThieleCompleteScaled.v), shows that class is this one with
its prices multiplied by $c$ (\code{scale\_complete\_at} and \code{unscale\_complete}), and shows the
clock fails at every unit (\code{clock\_not\_thiele\_complete\_at\_any\_unit}).

\subsection{The universal program's flag and ledger}

{\raggedright \code{universal\_halting}, \code{universal\_flag\_iff}, \code{universal\_ledger\_exact} and \code{universal\_thiele\_complete}, coq/kernel/foundation/UniversalRun.v.\par}

{\scriptsize
\begin{verbatim}
Theorem universal_halting : forall P x y,
  (exists m, E.halted P (E.core_of (grun P x y m))) <->
  (exists n, M.halted U (M.core_of (hrun P x y n))).

Theorem universal_flag_iff : forall P x y,
  (exists m, E.cert (grun P x y m) = true) <->
    (exists n, M.cert (hrun P x y n) = true).

Theorem universal_ledger_exact : forall P x y,
  (forall m, exists n,
     ((rel P (grun P x y m) (hrun P x y n) /\ m <= n) \/
      rel_halt P (grun P x y m) (hrun P x y n)) /\
     M.mu (hrun P x y n) = E.mu (grun P x y m)) /\
  (forall n, exists m,
     E.mu (grun P x y m) <= M.mu (hrun P x y n) <= E.mu (grun P x y (S m))).

Theorem universal_thiele_complete : TC.thiele_complete host_machine.
\end{verbatim}}

In plain words: \texttt{grun P x y m} is $m$ program steps of the small machine's program
$P$ from \texttt{start x y}; \texttt{hrun P x y n} is $n$ program steps of the one fixed host
program \texttt{U} from the host loaded with $P$, $x$ and $y$. The first theorem: the
guest halts at some step exactly when \texttt{U} does. The second: the guest's flag is up at some
step exactly when the host's is. The third: for every guest step count there's a host step count
where the two are related,
with the host count at least the guest count, or both stopped, with equal ledgers; and every host point's ledger lies
between the guest's at some step and at the step after. \texttt{rel} and \texttt{rel\_halt} are the
relations Theorem~\ref{thm:u34-U-sim} spells out. The fourth: the host machine, the small machine
with a register for every whole number and the one property \code{PSlot}, read one move per
instruction, is certified Thiele-complete in the sense of the previous subsection.
Theorem~\ref{thm:u34-U-flag} says exactly the second and the third. Theorem~\ref{thm:u34-U-halting}
says the first and more: at halting the registers, the traps, the ledgers and the flags match too
(\code{universal\_output}, not printed here). Theorem~\ref{thm:u34-U-complete} says the fourth
and more: $U$'s run is a run of the host, the floor on the cost of a run that raises the flag, and the
undecidability of $U$'s halting, each with its own Coq theorem in the table of
Section~\ref{sec:crosswalk}.

\subsection{The universal reading checker}

{\raggedright \code{cg\_run\_check}, \code{cg\_ueval} and \code{cg\_ueval\_iff}, coq/kernel/foundation/CompilerChecker.v.\par}

{\scriptsize
\begin{verbatim}
Definition cg_run_check (rc : cg_routine_code) (v : nat) : bool :=
  match rc with
  | (ig, R, xS, xB, xT, m) =>
      let st := cg_mme_run_fuel (cg_expo (qs xT) v) (ig, R) (ig, cg_env_chk m v) in
      cg_out_codeb (ig, R) (fst st) && Nat.eqb (get_env (snd st) xB) 1
  end.

Definition cg_ueval (p : cg_uprop) (v : nat) : bool :=
  match p with
  | UBase q => G.ceval q v
  | URun r =>
      match cg_rdec r with
      | Some rc => cg_run_check rc v
      | None => false
      end
  end.

Theorem cg_ueval_iff : forall p v, cg_ueval p v = true <-> cg_uholds p v.
\end{verbatim}}

In plain words: a claim is either one of the counter language's words, checked as before, or a
number $r$ naming a reading routine. For the second kind the checker decodes $r$; a code that
names no routine gets false. Otherwise it reads the registers and a fuel number $t$ out of the
value, runs the routine for at most $t$ steps, and says yes exactly when the routine has left its
code with 1 in its answer register. The theorem says the checker's yes is the claim's meaning,
which for a routine is that bounded run. Section~\ref{sec:presented-universal} says exactly this.
The checker answers on every input because decoding terminates and the routine runs for
at most the supplied finite number of steps
(\code{cg\_ueval\_total}).

\subsection{The surcharge can't always be avoided}

{\raggedright \code{presented\_universal\_no\_exact\_below\_three}, coq/kernel/foundation/PresentedUniversal.v.\par}

{\scriptsize
\begin{verbatim}
Theorem presented_universal_no_exact_below_three :
  forall (Mp : presented_machine) (s0 : T.cs_state (pm_sys Mp)) (n : nat),
  mlatch Mp s0 n = true -> mledger Mp s0 n < 3 ->
  forall (h0 : hstate) (tr : list (@M.pu_instr pu_hprop)),
  M.pu_clean_start h0 ->
  ~ (M.cert (hrun_l tr h0) = mlatch Mp s0 n /\
     M.mu (hrun_l tr h0) = M.mu h0 + mledger Mp s0 n).
\end{verbatim}}

In plain words: if a presented machine's latch is up after $n$ moves while its ledger is below
three, then no list of host instructions, run from a clean host start, ends with the host's flag
equal to that latch and the host's ledger raised by exactly that ledger.
Theorem~\ref{thm:u34-presented-surcharge}(3) says exactly this. It's about the guest's whole
ledger at that point. A guest whose one raising move is cheap but whose earlier moves already
paid three or more isn't covered.

\subsection{Nesting pays for every rise, on every host run that follows the guest}

{\raggedright \code{cmpz\_hom}, coq/kernel/foundation/CzCat.v, and \code{cmpz\_follows},
\code{cmpz\_follow\_cost\_ge\_exits} and \code{cmpz\_end\_only\_not\_enough},
coq/kernel/foundation/CzFollow.v; \code{cmpz\_follows} and \code{cmpz\_follow\_cost\_ge\_exits}
are written inside a section that fixes \texttt{X}, \texttt{Y} and \texttt{f : cmpz\_hom X Y}.\par}

{\scriptsize
\begin{verbatim}
Record cmpz_hom (X Y : cmpz_obj) : Type := mk_hom {
  h_state : am_state (ob_M X) -> am_state (ob_M Y);
  h_rec : ob_A X -> ob_A Y;
  h_mono : forall a b, bp_le (ob_P X) a b -> bp_le (ob_P Y) (h_rec a) (h_rec b);
  h_rec_nat : forall s, am_rec (ob_M Y) (h_state s) = h_rec (am_rec (ob_M X) s);
  h_sim : forall s m, exists tr : list (am_move (ob_M Y)),
            h_state (am_step (ob_M X) s m) = am_run (ob_M Y) tr (h_state s)
}.

Fixpoint cmpz_follows (tr : list (am_move (ob_M X))) (s : am_state (ob_M X))
    (segs : list (list (am_move (ob_M Y)))) : Prop :=
  match tr, segs with
  | [], [] => True
  | m :: tr', seg :: segs' =>
      h_state f (am_step (ob_M X) s m) = am_run (ob_M Y) seg (h_state f s) /\
      cmpz_follows tr' (am_step (ob_M X) s m) segs'
  | _, _ => False
  end.

Theorem cmpz_follow_cost_ge_exits :
  (forall a b, bp_le (ob_P Y) (h_rec f a) (h_rec f b) -> bp_le (ob_P X) a b) ->
  ax_a2 (X := am_axsys (ob_M Y)) ->
  forall tr s segs, cmpz_follows tr s segs ->
    cmpz_cost (ob_M Y) (concat segs) >= cmpz_exit_count (ob_M X) tr s.

Theorem cmpz_end_only_not_enough :
  exists (X Y : cmpz_obj) (f : cmpz_hom X Y),
    (forall a b, bp_le (ob_P Y) (h_rec f a) (h_rec f b) -> bp_le (ob_P X) a b) /\
    ax_a2 (X := am_axsys (ob_M Y)) /\
    exists (tr : list (am_move (ob_M X))) s (tr' : list (am_move (ob_M Y))),
      h_state f (am_run (ob_M X) tr s) = am_run (ob_M Y) tr' (h_state f s) /\
      cmpz_cost (ob_M Y) tr' < cmpz_exit_count (ob_M X) tr s.
\end{verbatim}}

In plain words: a morphism $f$ from a guest $X$ to a host $Y$ is a map of states and a monotone
map of records such that the record of an image state is the image of the record, and every move
of the guest is simulated by some list of the host's moves landing on the image of the state
after the move. \code{cmpz\_follows tr s segs} says the host's run is split into one stretch per
guest move, each stretch carrying the image of the state before that move to the image of the
state after it. The first theorem: if the record map reflects the order and $Y$ obeys its toll,
then every host run that follows a guest run in that sense costs at least the number of guest
steps whose new record isn't at or below the old one. The second: for host runs that only reach
the image of the guest's end state, no such bound holds in general. On a one-bit machine whose
move $b$ sets the bit to $b$, the guest run ``up, then down'' has one such step, and the empty host
run reaches the image of its end for nothing. Theorem~\ref{thm:parts-nesting}(1) says exactly
these, with ``faithfully'' meaning that $f$ is a morphism, and with \code{cmpz\_follow\_exists} and
\code{cmpz\_follow\_reaches} (CzFollow.v) for the existence of a following run and where it ends.
The earlier form, \code{cmpz\_sim\_cost\_ge\_exits} (CzCat.v), finds one host run reaching the end
image at that cost. It bounds the host's cost by the guest's rises. It says nothing about the
guest's other charges.

\subsection{The copy is the record, on the simpler host}

{\raggedright \code{record\_agreement}, minimal/UniversalThiele.v.\par}

{\footnotesize
\begin{verbatim}
Definition agrees (h : hstate) : Prop := mrec h = E.cert (gst h).

Theorem record_agreement : forall tr h, agrees h -> agrees (hrun tr h).
\end{verbatim}}

In plain words: the host keeps a mirror bit \code{mrec} beside the guest's state \code{gst},
and \texttt{agrees h} says the mirror equals the guest's own flag. The theorem: a host state that
agrees still agrees after any list of host moves, whatever they are.
Observation~\ref{thm:u34-mirror} says this for lists of moves and adds single moves, stored programs, loading,
and the mirror after $n$ moves GSTEP 1,
which are \code{record\_agreement\_step}, \code{record\_agreement\_program}, \code{hload\_agrees} and \code{simulated\_record}.

\subsection{No unpaid host certification, on the simpler host}

{\raggedright \code{no\_free\_host\_certification}, minimal/UniversalThiele.v, with the host's moves and their costs.\par}

{\scriptsize
\begin{verbatim}
Inductive hinstr : Type :=
| OWN (i : E.instr)
| GSTEP (b : nat).

Definition hcost (i : hinstr) : nat :=
  match i with OWN j => E.cost j | GSTEP b => b end.

Definition crossing (g g' : E.state) : bool := negb (E.cert g) && E.cert g'.

Theorem no_free_host_certification : forall tr h,
  mrec h = false -> mrec (hrun tr h) = true ->
  exists pre b post, tr = pre ++ GSTEP b :: post /\ b >= 1 /\
    crossing (gst (hrun pre h)) (gst (hexec (hrun pre h) (GSTEP b))) = true.
\end{verbatim}}

In plain words: a host move is either \texttt{OWN i}, one of the small machine's own
instructions on the host's core, or \texttt{GSTEP b}, one guest step with budget \code{b}, which
costs \code{b} whatever the guest instruction cost. \texttt{crossing g g'} says the guest's flag
went from down to up. The theorem: if the mirror is down before a list of host moves and up
after it, then somewhere in the list there's a \texttt{GSTEP b} with \code{b} at least one at
which the guest's flag crossed. Theorem~\ref{thm:u34-no-free-host} prints the one-move form,
\code{no\_free\_host\_certification\_step}, which also says the host wasn't trapped and the
guest's next instruction was CERTIFY; this is the form for whole lists that it lifts to. The printed theorem says two things more, each with its own Coq theorem: a trace of OWN moves alone leaves the mirror and
the guest state where they were (\code{guestless\_runs\_leave\_mirror}), and the host's own flag rises only by an OWN CERTIFY that passes (\code{own\_record\_only\_by\_certify}). Neither the one-move form nor the list form
says what the guest did before its CERTIFY. That, after a run from a loaded host, a raised mirror means
the guest's own run
checked and committed first is the next theorem.

\subsection{A raised mirror stands on the guest's own earned chain, on the simpler host}

{\raggedright \code{host\_mirror\_earned}, minimal/UniversalThiele.v.\par}

{\tiny
\begin{verbatim}
Theorem host_mirror_earned : forall tr a b P x y,
  mrec (hrun tr (hload a b P x y)) = true ->
  exists n gtr pre1 p c mid1 mid2 post,
    gtr = E.trace_of n P (E.start x y) /\
    gst (hrun tr (hload a b P x y)) = E.run gtr (E.start x y) /\
    gtr = pre1 ++ E.CHECK p c :: mid1 ++ E.COMMIT p c :: mid2 ++ E.CERTIFY :: post /\
    E.check_ok (E.core_of (E.run pre1 (E.start x y))) p c = true /\
    E.commit_ok (E.core_of (E.run (pre1 ++ E.CHECK p c :: mid1) (E.start x y))) p c
      = true /\
    E.certify_ok (E.core_of (E.run (pre1 ++ E.CHECK p c :: mid1 ++ E.COMMIT p c :: mid2)
                               (E.start x y))) = true /\
    E.ver (E.core_of (E.run pre1 (E.start x y))) c
      = E.ver (E.core_of (E.run (pre1 ++ E.CHECK p c :: mid1) (E.start x y))) c /\
    E.untouched (E.run (pre1 ++ [E.CHECK p c]) (E.start x y)) mid1 c.
\end{verbatim}}

In plain words: load the host with a guest program $P$ and a guest start \texttt{(x, y)}
(\code{hload} gives the guest a clean start), run any list of host moves, and suppose the
mirror is up at the end. Then for some number of steps \code{n}, the guest's state is its own
program's run of \code{n} steps, and the instructions that run executed split exactly the way
Theorem~\ref{thm:u34-cert-earned} splits a trace: a passing CHECK, a COMMIT of the same claim at
the same version with the counter untouched between, then a passing CERTIFY. So the host's
mirror can't be up unless the guest's own program earned it. The book's sentences citing
\code{host\_mirror\_earned} say this in words; the Coq says a little more than they need, since the chain
sits inside the instructions the guest program itself executed. Any list containing the chain
would have been enough for them.

\subsection{Structural entitlement, represented}

{\raggedright \code{ent\_representation}, minimal/EntitlementSmall.v, with the definitions it leans on.\par}

{\tiny
\begin{verbatim}
Definition ent_strict_sublist (post prior : list X) : Prop :=
  incl post prior /\ exists x, In x prior /\ ~ In x post.

Definition ent_distinguishes (obs : X -> O) (w : X) (post : list X) : Prop :=
  forall t, In t post -> obs w <> obs t.

Definition ent_member (eqb : O -> O -> bool) (obs : X -> O) (L : list X) (o : O) : bool :=
  existsb (fun t => eqb (obs t) o) L.

Definition ent_stronger (P Q : O -> bool) : Prop :=
  forall o, P o = true -> Q o = true.

Definition ent_strictly_stronger (P Q : O -> bool) : Prop :=
  ent_stronger P Q /\ exists o, P o = false /\ Q o = true.

Definition ent_reduction (r : X -> O) (T : ent_tree) (prior post : list X) : Prop :=
  exists F : X -> list X,
    (forall x, In x prior -> exists t, In t post /\ In x (F t) /\ r x = r t) /\
    length prior <= fold_right Nat.add 0 (map (fun t => length (F t)) post) /\
    (forall t, In t post -> length (F t) <= ent_leaves T).

Definition ent_certified {M : machine} (I : thiele_interface M) {O : Type}
    (s0 : m_state M) (tr : list (m_move M))
    (d : list (m_move M) -> O) (P : O -> bool) : Prop :=
  m_record M (run M tr s0) = true /\ P (d tr) = true.

Theorem ent_representation :
  forall (M : machine) (I : thiele_interface M) (X O : Type)
         (s0 : m_state M) (tr : list (m_move M))
         (d : list (m_move M) -> O) (obs r : X -> O) (eqb : O -> O -> bool)
         (prior post : list X) (T : ent_tree),
    thiele_complete_with I ->
    (forall o1 o2, eqb o1 o2 = true <-> o1 = o2) ->
    ent_strict_sublist post prior ->
    (exists w, In w prior /\ ~ In w post /\ ent_distinguishes obs w post) ->
    ti_clean I s0 ->
    ent_certified I s0 tr d (ent_member eqb obs post) ->
    ent_depth T <= record_moves I tr ->
    0 < length post ->
    ent_reduction r T prior post ->
    ent_strictly_stronger (ent_member eqb obs post) (ent_member eqb obs prior) /\
    earned_chain I s0 tr /\
    (exists pre c chk mid1 cmt rest,
       tr = pre ++ chk :: mid1 ++ cmt :: rest /\
       ti_kind I chk = KCheck c /\ ti_kind I cmt = KCommit c /\
       ti_meaning I c (run M pre s0) /\ ti_meaning I c (run M (pre ++ chk :: mid1) s0)) /\
    Nat.log2_up (length prior) - Nat.log2_up (length post) <= record_moves I tr /\
    record_moves I tr = ti_ledger I (run M tr s0) - ti_ledger I s0 /\
    ti_ledger I s0 + 3 <= ti_ledger I (run M tr s0).
\end{verbatim}}

In plain words: take a machine and an interface through which it is certified Thiele-complete, any type of
rivals $X$ and of observations $O$ with an exact equality test, a run $t$ from a clean
state $s_0$, and a decoder $d$ that reads the run's receipts, its list of moves. Hypotheses
in order: the posterior list is inside the prior and drops at least one rival; some dropped rival
shows an observation no survivor shows; the start is clean; the run ends with the record up and the
posterior's membership test accepts the decoded run; the tree is no deeper than the number of record
moves in the run; the posterior is non-empty; and the covering is by fibres, as
\code{ent\_reduction} spells out. Conclusions in order: the posterior's test is strictly stronger; the
run contains the earned chain; some CHECK of a claim and a later COMMIT of the same claim appear in the run, with the claim holding at both; the index
bits are at most the record moves; the record moves are exactly the ledger's rise; and the rise is at
least 3. Theorem~\ref{thm:ent-machine} says this, with the hypotheses numbered (M1) to (M6)
and the conclusions numbered (1) to (5), and it says one thing more in (2): there the meanings hold at the chain's own CHECK and COMMIT. The Coq's third conclusion is a separate existential that doesn't name
the chain; the printed form follows from it with the last two parts of clause (b), a passing check means the claim holds and an unchanged thing keeps it holding.

{\raggedright And its form with the narrowing made by the run, \code{ent\_run\_entitlement}, same file, with the definitions it leans on.\par}

{\tiny
\begin{verbatim}
Fixpoint ent_narrow (fs : list (X -> bool)) (L : list X) : list X :=
  match fs with [] => L | f :: fs' => ent_narrow fs' (filter f L) end.

Fixpoint ent_stages (fs : list (X -> bool)) (L : list X) : list (nat * nat) :=
  match fs with
  | [] => []
  | f :: fs' => (length L, length (filter f L)) :: ent_stages fs' (filter f L)
  end.

Definition ent_prod (l : list nat) : nat := fold_right Nat.mul 1 l.

Fixpoint ent_passes {M : machine} (I : thiele_interface M) (s : m_state M)
    (done tr : list (m_move M)) : list (list (m_move M) * ti_claim I) :=
  match tr with
  | [] => []
  | m :: rest =>
      let tail := ent_passes I (m_step M s m) (done ++ [m]) rest in
      match ti_kind I m with
      | KCheck c => if ti_check I s c then (done, c) :: tail else tail
      | _ => tail
      end
  end.

Definition ent_run_tests {M : machine} (I : thiele_interface M) (s0 : m_state M)
    (tr : list (m_move M)) : list (m_state M -> bool) :=
  map (fun pc x => ti_check I (run M (fst pc) x) (snd pc)) (ent_passes I s0 [] tr).

Definition ent_run_posterior {M : machine} (I : thiele_interface M) (s0 : m_state M)
    (tr : list (m_move M)) (prior : list (m_state M)) : list (m_state M) :=
  ent_narrow (ent_run_tests I s0 tr) prior.

Theorem ent_run_entitlement :
  forall (M : machine) (I : thiele_interface M),
    thiele_complete_with I ->
    forall (s0 : m_state M) (tr : list (m_move M)) (prior : list (m_state M)) (k : nat),
    ti_clean I s0 ->
    m_record M (run M tr s0) = true ->
    In s0 prior ->
    Forall (fun st => fst st <= 2 ^ k * snd st) (ent_stages (ent_run_tests I s0 tr) prior) ->
    incl (ent_run_posterior I s0 tr prior) prior /\
    In s0 (ent_run_posterior I s0 tr prior) /\
    (forall x, In x prior -> ~ In x (ent_run_posterior I s0 tr prior) ->
       exists pre c m rest, tr = pre ++ m :: rest /\ ti_kind I m = KCheck c /\
         ti_check I (run M pre s0) c = true /\ ti_check I (run M pre x) c = false) /\
    length prior * ent_prod (map snd (ent_stages (ent_run_tests I s0 tr) prior))
      = length (ent_run_posterior I s0 tr prior)
          * ent_prod (map fst (ent_stages (ent_run_tests I s0 tr) prior)) /\
    Nat.log2_up (length prior) - Nat.log2_up (length (ent_run_posterior I s0 tr prior))
      <= k * length (ent_passes I s0 [] tr) /\
    length (ent_passes I s0 [] tr) <= record_moves I tr /\
    record_moves I tr = ti_ledger I (run M tr s0) - ti_ledger I s0 /\
    Nat.log2_up (length prior) - Nat.log2_up (length (ent_run_posterior I s0 tr prior))
      <= k * (ti_ledger I (run M tr s0) - ti_ledger I s0) /\
    earned_chain I s0 tr /\
    ti_ledger I s0 + 3 <= ti_ledger I (run M tr s0).
\end{verbatim}}

In plain words: \code{ent\_passes} walks the run and keeps, for each CHECK that passes at the state
the run is in, the moves before it and its claim; \code{ent\_run\_tests} turns each into a test of
a candidate start $x$, replay those moves from $x$ and check the same claim; the posterior keeps
the candidates that pass every test, in order; and \code{ent\_stages} records the number of
candidates standing before and after each test. Hypotheses: the machine is certified Thiele-complete through
\texttt{I}, the start is clean, the run ends with the record up, the start is among the
candidates, and every test keeps at least a $2^{-k}$ share. Conclusions in order: the posterior is
inside the prior and keeps the start; every dropped candidate fails, replayed, a CHECK the run made
and passed; the exact product count; the index bits are at most $k$ times the tests passed; the
tests passed are at most the record moves; the record moves are the ledger's rise; the index bits
are at most $k$ times that rise; the earned chain; and a rise of at least 3.
Theorem~\ref{thm:ent-run-machine} (1) to (5) says exactly this; its part (6) is
\code{ent\_run\_round\_bits}, not printed here.

\subsection{A raised flag stands on the CHSH check}

{\raggedright \code{small\_chsh\_flag\_implies\_tsirelson}, coq/kernel/quantum/SmallChshMachine.v, over the host machine with the one property \code{small\_chsh\_PCHSH}.\par}

{\scriptsize
\begin{verbatim}
Theorem small_chsh_flag_implies_tsirelson :
  forall (s0 : @state small_chsh_prop) (tr : list (@instr small_chsh_prop)),
  clean_start s0 ->
  cert (run small_chsh_prop_eqb small_chsh_eval tr s0) = true ->
  exists pre1 c mid1 mid2 post,
    tr = pre1 ++ CHECK small_chsh_PCHSH c :: mid1
           ++ COMMIT small_chsh_PCHSH c :: mid2 ++ CERTIFY :: post /\
    let v := vals (core_of (run small_chsh_prop_eqb small_chsh_eval pre1 s0)) c in
    let t := small_chsh_tally_of v in
    small_chsh_check t = true /\
    small_chsh_meaning t /\
    (small_chsh_score t * small_chsh_score t <= 8)%R /\
    (Rabs (small_chsh_score t) <= 2 * sqrt 2)%R /\
    vals (core_of (run small_chsh_prop_eqb small_chsh_eval
                    (pre1 ++ CHECK small_chsh_PCHSH c :: mid1) s0)) c = v.
\end{verbatim}}

In plain words: start from a clean start and run any list of instructions of the host machine whose
only property is the CHSH one, and suppose the flag is up at the end. Then the list splits as some
instructions, a CHECK of the property at some counter $c$, some more, a COMMIT of the same
property at the same counter, some more, a CERTIFY, and the rest. Take the value $v$ that
counter $c$ held just before the CHECK, and read it as a tally of eight counts. That tally
passes the seven integer facts (\code{small\_chsh\_check}); it means what the check is for: every
question pair sampled and the pinned five-by-five symmetric and positive semidefinite
(\code{small\_chsh\_meaning}); its score $S$ has $S^2 \le 8$ and $|S| \le 2\sqrt 2$; and counter
$c$ still held $v$ when the COMMIT ran. Because that tally passes the check, its score is in fact
strictly below the bound, $S^2 < 8$ (\code{nec\_e\_check\_strict}), and passing tallies come as
close to 8 as you like (\code{nec\_e\_check\_sharp}), both in NecEChsh.v. Theorem~\ref{thm:chsh-earned} says exactly this in
its first part. It says more in two places that have their own Coq theorems, the stored-program form
and the converse that the three instructions certify exactly when the check passes.

\subsection{The recursion theorem up to the record, but not on every register}

{\raggedright \code{sm2\_kleene\_obs} and \code{sm2\_no\_exact}, coq/kernel/foundation/SmFixedPoint.v and coq/kernel/foundation/SmNoExact.v, with the relation they use. Here \code{M} is the host machine's module, \code{hends P x s} says program \code{P}, started with $x$ in register 1 and everything else clean, stops (a trap counts as stopping) in state \code{s}, and \code{sm\_computes\_map T F} says that on the number of every program \code{p}, \code{T} stops with the number of \code{F p}.\par}

{\scriptsize
\begin{verbatim}
Definition sm2_obs (s t : hstate) : Prop :=
  M.vals (M.core_of s) 0 = M.vals (M.core_of t) 0 /\
  M.err (M.core_of s) = M.err (M.core_of t) /\
  M.mu s = M.mu t /\ M.cert s = M.cert t /\
  length (M.facts (M.core_of s)) = length (M.facts (M.core_of t)) /\
  (M.chan (M.core_of s) = None <-> M.chan (M.core_of t) = None).

Definition sm2_obs_equiv (P Q : list hinstr) : Prop :=
  forall x,
    (forall s, hends P x s -> exists t, hends Q x t /\ sm2_obs s t) /\
    (forall t, hends Q x t -> exists s, hends P x s /\ sm2_obs s t).

Theorem sm2_kleene_obs : forall (F : list hinstr -> list hinstr) (T : list hinstr),
  sm_computes_map T F -> exists e, sm2_obs_equiv e (F e).

Theorem sm2_no_exact :
  exists (F : list hinstr -> list hinstr) (T : list hinstr),
    sm_computes_map T F /\ forall e, ~ hequiv e (F e).
\end{verbatim}}

In plain words: \code{sm2\_obs s t} says two final states agree on register 0, the trap latch, the
ledger, the flag, the number of facts and whether the channel is empty, and nothing else.
\code{sm2\_obs\_equiv P Q} says that on every input, any state \code{P} ends in has a state \code{Q}
ends in that agrees in that sense, and the same the other way round, which includes that \code{P}
stops on an input exactly when \code{Q} does. The first theorem: for every transformation $F$
of host programs that some host program \code{T} carries out on numbers, there is a program \code{e}
that agrees with \code{F e} in that sense. \code{T} is any host program, and may itself check, commit
and certify. The second: there is a transformation $F$, carried out by a host program, such that
no program \code{e} behaves the same as \code{F e} in the sense of Definition~\ref{def:host-same}, which
compares every register, the fact table, the channel, the trap latch, the ledger and the flag.
Theorems~\ref{thm:host-kleene-record} and~\ref{thm:host-no-exact} say these two. The book's printed
version says exactly this, with the line number and versions left out of both comparisons, as the
definition does.

%===========================================================================
